\section{Test Scenarios}
\label{sec:scenarios}

We developed three comprehensive test scenarios to validate our distributed MPI WiFi implementation across different use cases: basic integration testing, realistic home WiFi networks, and large-scale vehicular communication.

\subsection{Phase 2.2.3c Integration Test}

\subsubsection{Purpose and Scope}

This fundamental test validates bidirectional MPI communication between channel and device processors. It serves as a minimal working example demonstrating:
\begin{itemize}
    \item Device registration with the channel processor
    \item Packet transmission via MPI messages
    \item RX notification delivery back to devices
    \item Multi-machine execution verification
\end{itemize}

\subsubsection{Topology}

\begin{itemize}
    \item \textbf{Ranks}: 4 total (1 channel + 3 device groups)
    \item \textbf{Devices}: 6 total (2 per device rank)
    \item \textbf{Positions}: Three groups at 0m, 50m, and 100m
    \item \textbf{Mobility}: Static (ConstantPositionMobilityModel)
    \item \textbf{Propagation}: Log-distance model (exponent 2.5)
\end{itemize}

\subsubsection{Traffic Generation}

Simple UDP applications with varied parameters:
\begin{lstlisting}[style=cpp, caption={Integration test traffic configuration}]
// Different TX powers per device group
Rank 1 devices: 20 dBm, 15 dBm
Rank 2 devices: 10 dBm, 12 dBm  
Rank 3 devices: 25 dBm, 18 dBm

// UDP echo applications
Port: 4444
Packet size: 300 bytes
Interval: 1 second
Duration: 10 seconds
\end{lstlisting}

\subsubsection{Validation Criteria}

The test verifies:
\begin{enumerate}
    \item All devices successfully register with rank 0
    \item TX\_REQUEST messages reach channel processor
    \item Propagation calculations complete
    \item RX\_NOTIFICATION messages return to device ranks
    \item Correct hostname/PID reporting for multi-machine execution
\end{enumerate}

\subsection{Home WiFi Scenario}

\subsubsection{Motivation}

Modern homes contain diverse WiFi devices with vastly different communication patterns. This scenario models a realistic residential deployment to test:
\begin{itemize}
    \item Heterogeneous device types and traffic loads
    \item Indoor propagation characteristics
    \item Various TX power levels (3-25 dBm)
    \item Sustained high data rates (50+ Mbps)
\end{itemize}

\subsubsection{Device Types and Specifications}

Table~\ref{tab:home-wifi-devices} details the 8 simulated devices:

\begin{table}[htbp]
\centering
\caption{Home WiFi device specifications}
\label{tab:home-wifi-devices}
\small
\begin{tabular}{@{}llcccc@{}}
\toprule
\textbf{Device} & \textbf{Type} & \textbf{Position (m)} & \textbf{TX Power} & \textbf{Data Rate} & \textbf{Pattern} \\
\midrule
Living Room TV & Smart TV & (5, 3, 1.0) & 15 dBm & 50 Mbps & Streaming \\
Kitchen Tablet & Tablet & (8, 8, 1.2) & 10 dBm & 10 Mbps & Browsing \\
Bedroom Laptop & Laptop & (12, 5, 1.0) & 20 dBm & 50 Mbps & Streaming \\
Kids Room Phone & Smartphone & (12, 8, 1.0) & 8 dBm & 10 Mbps & Social media \\
Security Camera & IoT & (2, 10, 2.5) & 12 dBm & 25 Mbps & Video upload \\
Smart Speaker & IoT & (6, 6, 1.0) & 5 dBm & 1 Mbps & Voice \\
Gaming Console & Console & (4, 4, 0.8) & 18 dBm & 30 Mbps & Gaming bursts \\
Smart Thermostat & IoT & (9, 2, 1.5) & 3 dBm & 1 Mbps & Periodic \\
\bottomrule
\end{tabular}
\end{table}

\subsubsection{Topology and Layout}

The scenario models a 15m × 12m home with devices positioned across different rooms. Figure~\ref{fig:home-layout} shows the floor plan:

\begin{figure}[htbp]
\centering
\begin{verbatim}
                     15 meters
        ┌───────────────────────────────────┐
        │                                   │
  12m   │  [Sec Cam]                        │
        │   (2,10)        Kitchen           │
        │              [Tablet] (8,8)       │
        │                                   │
        │                        [Phone]    │
        │  [Speaker]              (12,8)    │
        │   (6,6)      Living                │
        │               Room     Bedroom    │
        │  [Console]                        │
        │   (4,4)                           │
        │                       [Laptop]    │
        │  [TV]                  (12,5)     │
        │   (5,3)                           │
        │                                   │
        │           [Thermostat]            │
        │             (9,2)                 │
        └───────────────────────────────────┘
\end{verbatim}
\caption{Home WiFi device layout (top-down view)}
\label{fig:home-layout}
\end{figure}

\subsubsection{Propagation Environment}

Indoor propagation uses a log-distance model with parameters reflecting residential environments:
\begin{itemize}
    \item \textbf{Path loss exponent}: 3.0 (typical for indoor with walls)
    \item \textbf{Reference distance}: 1.0 meter
    \item \textbf{Reference loss}: 40 dB (2.4 GHz indoor)
\end{itemize}

At these settings, devices 10 meters apart experience ~70 dB path loss, suitable for testing realistic SNR conditions.

\subsubsection{Traffic Patterns}

Devices generate traffic matching their real-world counterparts:

\paragraph{Heavy Streamers (TV, Laptop)}
\begin{lstlisting}[style=cpp]
DataRate: 50 Mbps (OnOff, constant)
PacketSize: 1400 bytes
Application: Video streaming simulation
\end{lstlisting}

\paragraph{Medium Browsers (Phone, Tablet)}
\begin{lstlisting}[style=cpp]
DataRate: 10 Mbps (OnOff, constant)
PacketSize: 1200 bytes
Application: Web browsing, social media
\end{lstlisting}

\paragraph{Constant Uploaders (Security Camera)}
\begin{lstlisting}[style=cpp]
DataRate: 25 Mbps (OnOff, constant)
PacketSize: 1400 bytes
Application: HD video upload
\end{lstlisting}

\paragraph{Bursty Traffic (Gaming Console)}
\begin{lstlisting}[style=cpp]
DataRate: 30 Mbps peak
OnTime: Exponential(mean=2s)
OffTime: Exponential(mean=0.5s)
PacketSize: 600 bytes
Application: Gaming with bursts
\end{lstlisting}

\paragraph{Light IoT (Speaker, Thermostat)}
\begin{lstlisting}[style=cpp]
DataRate: 1 Mbps (OnOff, constant)
PacketSize: 200-400 bytes
Application: Voice commands, sensor data
\end{lstlisting}

\subsubsection{MPI Distribution}

Devices are distributed across 4 MPI ranks:
\begin{itemize}
    \item \textbf{Rank 0}: Channel processor (no devices)
    \item \textbf{Rank 1}: Living Room TV, Kitchen Tablet
    \item \textbf{Rank 2}: Bedroom Laptop, Kids Room Phone
    \item \textbf{Rank 3}: Security Camera, Smart Speaker, Gaming Console, Smart Thermostat
\end{itemize}

This distribution balances computational load while ensuring cross-rank communication (devices in different rooms often communicate through the shared channel).

\subsubsection{Duration and Scale}

\begin{itemize}
    \item \textbf{Simulation time}: 300 seconds (5 minutes)
    \item \textbf{Warm-up period}: 2 seconds (stagger application starts)
    \item \textbf{Total expected packets}: ~1.8 million across all devices
    \item \textbf{Aggregate offered load}: ~177 Mbps
\end{itemize}

\subsection{Home WiFi with PCAP}

\subsubsection{Hybrid Architecture Application}

This variant applies our hybrid MPI + local monitoring approach (Section~\ref{sec:hybrid}) to the home WiFi scenario. Each device has:

\begin{enumerate}
    \item \textbf{MPI Interface}: Heavy traffic for performance testing
    \begin{itemize}
        \item Same specifications as base scenario
        \item Distributed via \texttt{RemoteYansWifiChannelStub}
        \item IP subnet: 192.168.x.0/24
    \end{itemize}
    
    \item \textbf{Local Monitoring Interface}: Light traffic for PCAP
    \begin{itemize}
        \item Reduced data rates (1-5 Mbps)
        \item Local \texttt{YansWifiChannel}
        \item PCAP capture enabled
        \item IP subnet: 10.x.0.0/24
    \end{itemize}
\end{enumerate}

\subsubsection{PCAP Capture Configuration}

\begin{lstlisting}[style=cpp, caption={PCAP enable on monitoring interfaces}]
// Enable promiscuous mode PCAP on local devices
localPhy.EnablePcapAll("home-wifi-rank" + 
                       std::to_string(systemId), 
                       true);  // Promiscuous mode

// Generates files like:
// home-wifi-rank1-0-1.pcap (TV monitoring interface)
// home-wifi-rank1-1-1.pcap (Tablet monitoring interface)
// home-wifi-rank2-0-1.pcap (Laptop monitoring interface)
// ...
\end{lstlisting}

\subsubsection{Traffic Distribution}

Table~\ref{tab:hybrid-traffic} compares MPI vs. monitoring traffic:

\begin{table}[htbp]
\centering
\caption{Traffic distribution in hybrid scenario}
\label{tab:hybrid-traffic}
\small
\begin{tabular}{@{}lccc@{}}
\toprule
\textbf{Device} & \textbf{MPI Rate} & \textbf{Monitor Rate} & \textbf{Purpose} \\
\midrule
Smart TV & 50 Mbps & 5 Mbps & Performance + Protocol \\
Laptop & 50 Mbps & 5 Mbps & Performance + Protocol \\
Tablet & 10 Mbps & 2 Mbps & Medium load + Protocol \\
Phone & 10 Mbps & 2 Mbps & Medium load + Protocol \\
Security Cam & 25 Mbps & 5 Mbps & Sustained + Protocol \\
Gaming Console & 30 Mbps peak & 5 Mbps & Bursts + Protocol \\
Smart Speaker & 1 Mbps & 1 Mbps & Light + Protocol \\
Thermostat & 1 Mbps & 1 Mbps & Periodic + Protocol \\
\midrule
\textbf{Total} & \textbf{177 Mbps} & \textbf{26 Mbps} & \\
\bottomrule
\end{tabular}
\end{table}

\subsection{V2X Vehicular Communication Scenario}

\subsubsection{Motivation}

Vehicular networks represent one of the most demanding WiFi use cases, requiring:
\begin{itemize}
    \item Large node counts (200-500+ vehicles in highway scenarios)
    \item High mobility (30-40 m/s vehicle speeds)
    \item Frequent beacon broadcasts (10-20 Hz safety messages)
    \item Heterogeneous vehicle types (cars, trucks, emergency, buses)
    \item Infrastructure elements (RSUs - Road Side Units)
\end{itemize}

This scenario validates our implementation's scalability for mission-critical applications.

\subsubsection{Highway Topology}

\paragraph{Physical Layout}
\begin{itemize}
    \item \textbf{Lanes}: 4 (two directions, two lanes each)
    \item \textbf{Lane width}: 3.5 meters (standard highway)
    \item \textbf{Vehicles per lane}: 50
    \item \textbf{Initial spacing}: 20 meters
    \item \textbf{Total highway length}: 1000 meters
    \item \textbf{RSUs}: 5 units, spaced 200m apart
\end{itemize}

\paragraph{Vehicle Distribution}
\begin{itemize}
    \item \textbf{Total vehicles}: 200
    \item \textbf{Passenger cars}: 70\% (every regular position)
    \item \textbf{Trucks}: 20\% (every 5th vehicle)
    \item \textbf{Emergency vehicles}: 5\% (every 20th vehicle)
    \item \textbf{Motorcycles}: 5\% (every 8th vehicle)
    \item \textbf{Buses}: 5\% (every 12th vehicle)
\end{itemize}

\subsubsection{Vehicle Types and Communication}

Table~\ref{tab:v2x-vehicles} specifies vehicle communication parameters:

\begin{table}[htbp]
\centering
\caption{V2X vehicle type specifications}
\label{tab:v2x-vehicles}
\small
\begin{tabular}{@{}lcccp{4cm}@{}}
\toprule
\textbf{Type} & \textbf{Beacon Rate} & \textbf{TX Power} & \textbf{Payload} & \textbf{Purpose} \\
\midrule
Passenger Car & 10 Hz & 20 dBm & 300 B & Safety beacons (CAM) \\
Truck & 8 Hz & 23 dBm & 600 B & Fleet telemetry + CAM \\
Emergency & 20 Hz & 25 dBm & 400 B & High-priority warnings \\
Motorcycle & 15 Hz & 18 dBm & 200 B & High-frequency safety \\
Bus & 5 Hz & 22 dBm & 500 B & Transit info + CAM \\
RSU & 1 Hz & 27 dBm & 800 B & Infrastructure beacons \\
\bottomrule
\end{tabular}
\end{table}

CAM = Cooperative Awareness Message (ETSI standard)

\subsubsection{Mobility Model}

Vehicles use \texttt{ConstantVelocityMobilityModel}:

\begin{lstlisting}[style=cpp, caption={V2X mobility configuration}]
// Base velocity: 30 m/s (108 km/h)
// Add random variation per vehicle
Ptr<UniformRandomVariable> speedVar = 
    CreateObject<UniformRandomVariable>();
double speed = 30.0 + speedVar->GetValue(-5.0, 5.0);

// Set velocity (movement in +X direction)
mobility->SetVelocity(Vector(speed, 0.0, 0.0));

// Initial positions: staggered by lane and spacing
double x = vehicleIndexInLane * 20.0;
double y = laneNumber * 3.5;
double z = 1.5;  // Vehicle antenna height
mobility->SetPosition(Vector(x, y, z));
\end{lstlisting}

Over 300 seconds, vehicles travel ~9km, ensuring significant topology changes.

\subsubsection{V2X Propagation}

Two-ray ground model matches vehicular environments:

\begin{lstlisting}[style=cpp, caption={V2X propagation configuration}]
Ptr<TwoRayGroundPropagationLossModel> twoRay = 
    CreateObject<TwoRayGroundPropagationLossModel>();

// Vehicle antenna height: 1.5m
// RSU antenna height: 5.0m
twoRay->SetAttribute("HeightAboveZ", DoubleValue(1.5));

// Frequency: 5.9 GHz (802.11p/DSRC band)
twoRay->SetAttribute("Frequency", DoubleValue(5.9e9));
\end{lstlisting}

This model accounts for ground reflections, critical for vehicular scenarios.

\subsubsection{MPI Distribution Strategy}

For 200 vehicles across 4-6 ranks:

\paragraph{4-Rank Configuration}
\begin{itemize}
    \item Rank 0: Channel processor + 5 RSUs
    \item Rank 1: Vehicles 0-66 (lanes 0-1, first half)
    \item Rank 2: Vehicles 67-133 (lanes 0-1, second half)
    \item Rank 3: Vehicles 134-199 (lanes 2-3, all)
\end{itemize}

\paragraph{6-Rank Configuration} (better load balance)
\begin{itemize}
    \item Rank 0: Channel processor + 5 RSUs
    \item Ranks 1-5: 40 vehicles each
\end{itemize}

\subsubsection{Traffic Load Analysis}

With 200 vehicles and varied beacon rates:
\begin{align}
\text{Total beacons/sec} &= 140 \times 10 + 40 \times 8 + 10 \times 20 \notag \\
&\quad + 10 \times 15 + 10 \times 5 + 5 \times 1 \notag \\
&= 1400 + 320 + 200 + 150 + 50 + 5 \notag \\
&= 2125 \text{ beacons/second}
\end{align}

Channel processor performs ~425K propagation calculations/second (2125 TX × 200 receivers).

\subsubsection{Simulation Parameters}

\begin{itemize}
    \item \textbf{Duration}: 300 seconds (5 minutes)
    \item \textbf{WiFi standard}: 802.11a (basis for 802.11p)
    \item \textbf{Data mode}: OfdmRate6Mbps (robust for mobile)
    \item \textbf{MAC protocol}: AdHoc (no association overhead)
    \item \textbf{Expected total packets}: ~638K across all vehicles
    \item \textbf{Aggregate beacon load}: ~6.4 Mbps
\end{itemize}

\subsection{Scenario Comparison}

Table~\ref{tab:scenario-comparison} summarizes key differences:

\begin{table}[htbp]
\centering
\caption{Test scenario comparison}
\label{tab:scenario-comparison}
\small
\begin{tabular}{@{}lccc@{}}
\toprule
\textbf{Aspect} & \textbf{Integration} & \textbf{Home WiFi} & \textbf{V2X} \\
\midrule
Purpose & Basic validation & Realistic load & Scalability \\
Devices & 6 & 8 & 200 \\
Mobility & Static & Static & Mobile (30 m/s) \\
Duration & 10 s & 300 s & 300 s \\
Ranks & 4 & 4 & 4-6 \\
Propagation & Log-distance & Log-distance & Two-ray ground \\
Traffic & Light & Heavy (177 Mbps) & Beacons (2K/s) \\
PCAP & No & Optional & Optional \\
Complexity & Low & Medium & High \\
\bottomrule
\end{tabular}
\end{table}

These scenarios provide comprehensive coverage from basic functionality through realistic deployment to extreme scale, validating our implementation across the spectrum of WiFi use cases.
